\section{A smooth coefficient domain for the original relative ample model}

\begin{lemma}[Smooth arithmetic coefficients retain the entire original ample line]
\label{lem:kn-relative-arithmetic-smooth-ample-model}
\lean{EnvelopeProperty.Vanishing.RelativeArithmeticSmoothAmpleModel.exists_original_smooth_ample_model}
\uses{lem:kn-relative-arithmetic-original-ample-model,lem:kn-relative-arithmetic-scalar-transition,def:arithmetic-actual-localized-coefficient-ring}
\leanok
The original relative ample arithmetic model can be chosen over
an actual integer subalgebra $A\subset k$ which is smooth and
of finite type over $\mathbb Z$.  It retains the separate
finitely presented affine-base algebra with original fibre $R$,
the total smooth dimension $n$, separatedness, the actual
relative ampleness, and the whole original-line pullback.
\end{lemma}
\begin{proof}
Take the produced ample model over its finite-type integer
coefficient domain $A_0\subset k$.  The generic point of $A_0$
lies in its smooth locus over $\mathbb Z$, so choose a nonzero
$s\in A_0$ such that $A_0[1/s]$ is smooth.  Its actual embedding
in $k$ identifies it with an integer subalgebra $A\subset k$
containing $A_0$.  This algebra is of finite type over
$\mathbb Z$ and, therefore, also of finite type over $A_0$.
Regard it as a finite coefficient stage and apply the actual
coefficient-transition lemma to the whole prior family.
Smoothness of total dimension $n$, separatedness, and
quasi-compactness are preserved by this base change, as is the
actual relative ampleness of the same model line.  The modeled
affine algebra becomes its actual tensor extension, remains
finitely presented, and has original field fibre $R$ by tensor
associativity.  The transition's original-field square and
whole-module pullback isomorphism retain the same original $L$.
\end{proof}
